% Judul penuh meluber 14,54-20,02 pt di kepala halaman; hanya mark header yang
% dipendekkan (sekali pakai, lihat \judulpendek di preamble-en.tex).
\judulpendek{Law of Large Numbers and Central Limit Theorem}
\chapter{The Law of Large Numbers and the Central Limit Theorem}

A short program computes the average of $\exp(U)$ over ten thousand random numbers
$U$ uniform on $(0,1)$. The result it prints is $1.722206$. The true value is
$e-1=1.7182818285$, so the error of that run is $+0.003924$. The program is run
again with four times as many numbers: $1.717502$, error $-0.000780$. With ten
times as many: $1.717767$, error $-0.000515$. The pattern is not tidy. Four
times the data cuts the error by more than a factor of five, whereas ten times the
data cuts it by only $7.6$ times. The naive guess "ten times the data, one tenth
of the error" does not hold for a single run. One run of a program carries no rate
at all; what has a rate is the distribution over many runs.

That distribution can be measured. With $n=10^4$ the standard deviation of this
estimator is $0.00492$; with $n=4\cdot10^4$ it becomes $0.00246$; with $n=10^5$
it becomes $0.00156$. Every time $n$ grows fourfold, that deviation is exactly
halved. Here we have two different statements that are often confused. First, the
average really does approach the true value: that is a tendency, and its name is
the law of large numbers. Second, the remaining error has a definite size and
shape: that is fluctuation, and its name is the central limit theorem. This chapter
supplies names, proofs, and limits of validity for both.

Chapter 18 has already prepared the ingredients: expectation, variance, and the
distribution of a sum through convolution. Convolution is useful only when $n$ is
small and the original shape is simple; for most distributions no closed form
exists. What we need is not the exact distribution at one $n$, but a statement that
holds for every large $n$. The practical questions are two: how close may the
average be trusted, and how large an error is reasonable. The answers come from
examining how the sequence of distributions behaves, not from computing one
distribution.

\begin{untukanak}
Imagine a class is asked to guess the number of marbles in a large jar. The
children's guesses miss badly in both directions. One child guesses three hundred,
another guesses two thousand, when the jar holds eight hundred. The teacher then
writes all the guesses on the board and computes their average. The number turns
out to be very close to eight hundred.

Now the teacher adds students. With thirty guessers the average still wobbles. With
three hundred guessers the average is calmer. With three thousand guessers the
average hardly shifts any more. There are two strange things here. First, not one
of the guessers has become better; only the number of guessers changed. Second,
adding guessers from thirty to three hundred helps far more than adding from three
hundred to three thousand. The help keeps growing, but more and more slowly.

The question that makes us curious: if each guess can miss by a great deal, why
can the average be trusted? And why does the help slow down in a pattern we can
predict, rather than at random? In this chapter we will find out why, and at the
same time find out when the habit of trusting the average deceives us.
\end{untukanak}

\section{Three Meanings of Approaching for a Random Variable}
\tingkat{basic}

In everyday language: a sequence of random variables can "approach" in three
different senses, and the loosest sense says nothing at all about the values coming
closer.

A sequence of real numbers has only one sense of approaching, namely that the
distance $|x_n-L|$ becomes small. For random variables that distance is itself a
random variable, so we have to choose which quantity is made small. The first
choice is the probability. This is the choice that is weakest intuitively but most
useful practically.

\kamus{convergence in probability}{konvergensi dalam peluang}{a promise that the
probability of missing by more than a tolerance shrinks to zero, not a promise that
the value comes closer on a single run}

\begin{definisi}[Convergence in probability]
A sequence of random variables $X_1,X_2,\dots$ \emph{converges in probability}
\index{convergence in probability} to a random variable $X$ if for every
$\varepsilon>0$
\[
P\bigl(|X_n-X|>\varepsilon\bigr)\to 0 \qquad (n\to\infty).
\]
\end{definisi}

Note where the quantifier sits. The value of $\varepsilon$ is chosen first, and the
probability may shrink only after that. This definition does not say that
$|X_n-X|$ shrinks for one particular outcome of the experiment; it says that the
event "missing by more than $\varepsilon$" becomes rarer and rarer. On the same
experiment, $X_n$ may still miss by a lot, provided that event is rare.

\kamus{convergence in distribution}{konvergensi dalam distribusi}{what resembles is
the shape of the distribution, not the value itself; two variables can share a shape
without ever coming close}

\begin{definisi}[Convergence in distribution]
Let $F_n$ be the distribution function of $X_n$ and $F$ the distribution function
of $X$. We write $X_n\rightsquigarrow X$ when $F_n(t)\to F(t)$ at every point $t$
at which $F$ is continuous. That limit is called the \emph{limit in
distribution}\index{convergence in
distribution}.
\end{definisi}

The condition "at continuity points" is not a technical whim. Take $X_n$ equal to
$1/n$ with probability one, and $X$ equal to $0$ with probability one. Then
$F_n(0)=0$ for every $n$, whereas $F(0)=1$, so the sequence does not converge at
$t=0$. At every $t\neq0$ the sequence does converge, and $t=0$ is the only point of
discontinuity of $F$. Without that condition, a correct conclusion would be
rejected because of a single point.

The third sense is stronger and more literal: almost every outcome of the
experiment, as one point $\omega$ of the sample space, produces a convergent
sequence of real numbers.

\begin{definisi}[Almost sure convergence]
$X_n$ \emph{converges almost surely}\index{convergence almost surely} to $X$ when
\[
P\bigl(\{\omega : X_n(\omega)\to X(\omega)\}\bigr)=1 .
\]
\end{definisi}

Here the sequence of real numbers $X_1(\omega),X_2(\omega),\dots$ really does
converge in the sense of Chapter 4, and the set of outcomes that fail to converge
may only have probability zero. Therefore almost sure convergence speaks about one
long run, whereas convergence in probability speaks about one $n$ across many runs
at once.

The three senses are not equivalent, but there is a fixed direction of implication.

\begin{teorema}[Relations among the three modes]
For random variables $X_n$ and $X$:
(a) $X_n\to X$ almost surely implies $X_n\to X$ in probability;
(b) $X_n\to X$ in probability implies $X_n\rightsquigarrow X$;
(c) if $X_n\rightsquigarrow c$ and $c$ is a constant, then $X_n\to c$ in
probability.
\end{teorema}

\begin{proof}[Proof]
For (a), fix $\varepsilon>0$ and let
$A_n=\{\sup_{k\ge n}|X_k-X|>\varepsilon\}$. The sets $A_n$ decrease, and every
$\omega$ in $\bigcap_n A_n$ satisfies $|X_k(\omega)-X(\omega)|>\varepsilon$ for
infinitely many $k$, so the sequence does not converge to $X(\omega)$. Hence
$\bigcap_n A_n$ is contained in a set of probability zero. Since probability is
continuous from above on sets of finite probability,
$P(A_n)\to P(\bigcap_n A_n)=0$. The event $\{|X_n-X|>\varepsilon\}$ is contained in
$A_n$, so its probability also tends to zero.

For (b), take $t$ a continuity point of $F$ and $\varepsilon>0$. From
$\{X_n\le t\}\subseteq\{X\le t+\varepsilon\}\cup\{|X_n-X|>\varepsilon\}$ we obtain
\[
F_n(t)\le F(t+\varepsilon)+P\bigl(|X_n-X|>\varepsilon\bigr).
\]
In the same way, from $\{X\le t-\varepsilon\}\subseteq\{X_n\le t\}\cup
\{|X_n-X|>\varepsilon\}$ we obtain
$F(t-\varepsilon)\le F_n(t)+P(|X_n-X|>\varepsilon)$. The two inequalities give
\[
F(t-\varepsilon)-\delta_n\le F_n(t)\le F(t+\varepsilon)+\delta_n,
\qquad \delta_n=P\bigl(|X_n-X|>\varepsilon\bigr).
\]
Since $\delta_n\to0$, the limit inferior and the limit superior of $F_n(t)$ are
trapped between $F(t-\varepsilon)$ and $F(t+\varepsilon)$. Let
$\varepsilon\downarrow0$ and use the continuity of $F$ at $t$; both tend to $F(t)$.

For (c), take $\varepsilon>0$. Since $F$ is the point-mass distribution at $c$, the
points $c-\varepsilon$ and $c+\varepsilon$ are both continuity points of $F$, so
\[
P\bigl(|X_n-c|>\varepsilon\bigr)
\le P(X_n<c-\varepsilon)+P(X_n>c+\varepsilon)
\le F_n(c-\varepsilon)+1-F_n(c+\varepsilon)
\to 0+1-1=0 .
\]
\end{proof}

The three reverse implications fail, and the counterexamples are worth remembering
because each one shows what is actually lost.

\begin{contoh}[In probability but not almost surely]
Split $(0,1)$ into $k$ pieces of equal length, for $k=1,2,3,\dots$, then arrange all
those pieces into one list. Define $X_n(\omega)=1$ if $\omega$ lies in the $n$-th
piece and $0$ otherwise. The $n$-th piece comes from the split into $k_n$ pieces, so
$P(X_n=1)=1/k_n\to0$ and $X_n\to0$ in probability. But for every
$\omega\in(0,1)$, within every block of splits there is exactly one $n$ with
$X_n(\omega)=1$. That value one appears infinitely often, so $X_n(\omega)$ does not
converge, for every $\omega$. Hence there is no almost sure convergence. This
sequence even converges in $L^1$, because $\mathbb{E}|X_n|=1/k_n\to0$, so $L^1$
strength does not close that gap either.
\end{contoh}

\begin{contoh}[In distribution but not in probability]
Take $X\sim N(0,1)$ and set $X_n=-X$ for every $n$. The distribution of $X_n$ is
exactly the distribution of $X$, so $F_n(t)=F(t)$ and $X_n\rightsquigarrow X$
vacuously. But $P(|X_n-X|>\varepsilon)=P(2|X|>\varepsilon)>0$ for every
$\varepsilon$, and that probability does not depend on $n$. Two variables of the
same shape need not be close: $X$ and $-X$ can be far apart on the same run.
\end{contoh}

The lesson of those two examples: convergence in distribution touches the
distributions one at a time, whereas the other two modes touch the relation between
variables on the same sample space. That is why part (c) of the theorem matters
practically: in asymptotic statistics the limit is almost always a constant, and
there the three modes meet.

One further warning has to be recorded from the start, because the mistake is
expensive.

\begin{catatan}
Convergence in probability does not carry convergence of expectations with it.
Suppose $P(X_n=n^2)=1/n$ and $P(X_n=0)=1-1/n$. For every $\varepsilon>0$,
$P(|X_n|>\varepsilon)=1/n\to0$, so $X_n\to0$ in probability. But
$\mathbb{E}[X_n]=n$, so the expectation tends to infinity. An average that
concentrates is no guarantee that its expectation concentrates.
\end{catatan}

\section{Chebyshev's Inequality and the Weak Law of Large Numbers}
\tingkat{intermediate}

In everyday language: the probability of straying far from the mean is bounded by
the variance divided by the square of that stray distance.

To prove the law of large numbers we need only one inequality, and that inequality
needs only one simple observation about expectation.

\begin{lemma}[Markov]
If $Y\ge0$ and $a>0$, then $P(Y\ge a)\le \mathbb{E}[Y]/a$.
\index{Markov's inequality}
\end{lemma}

\begin{proof}[Proof]
Write $Y=Y\mathbf{1}_{\{Y<a\}}+Y\mathbf{1}_{\{Y\ge a\}}$. The first term is
nonnegative. On the event $Y\ge a$ we have $Y\ge a$, so the expectation of the
second term is at least $a\,P(Y\ge a)$. Hence
$\mathbb{E}[Y]\ge 0+a\,P(Y\ge a)$.
\end{proof}

\begin{akibat}[Chebyshev's inequality]
If $\mathbb{E}[X]=\mu$ and $\operatorname{Var}(X)=\sigma^2<\infty$, then for every
$a>0$
\[
P\bigl(|X-\mu|\ge a\bigr)\le\frac{\sigma^2}{a^2}.
\]
\index{Chebyshev's inequality}
\end{akibat}

\begin{proof}[Proof]
Apply Markov's lemma to $Y=(X-\mu)^2$ with threshold $a^2$. The event
$\{Y\ge a^2\}$ is the same as $\{|X-\mu|\ge a\}$, and $\mathbb{E}[Y]=\sigma^2$.
\end{proof}

This inequality is weak, and its weakness is not a hidden defect but its strength.
It knows nothing about the shape of any distribution, only two moments. Therefore it
holds for every distribution with finite variance, and precisely for that reason it
is enough to prove a theorem that holds in general.

\kamus{law of large numbers}{hukum bilangan besar}{a promise that the average of
many observations sticks to the true mean, with a probability of missing that
shrinks}

\begin{teorema}[Weak law of large numbers]
Suppose $X_1,X_2,\dots$ are independent and identically distributed with
$\mathbb{E}[X_i]=\mu$ and $\operatorname{Var}(X_i)=\sigma^2<\infty$. Then for every
$\varepsilon>0$\index{law of large numbers}
\[
P\bigl(|\bar X_n-\mu|\ge\varepsilon\bigr)\to0,
\qquad \bar X_n=\frac1n\sum_{i=1}^n X_i .
\]
\end{teorema}

\begin{proof}[Proof]
Since expectation is linear, $\mathbb{E}[\bar X_n]=\mu$. Since the variables are
independent, the variance of a sum is the sum of the variances, so
$\operatorname{Var}(\bar X_n)=\sigma^2/n$. Chebyshev's inequality with
$a=\varepsilon$ gives
\[
P\bigl(|\bar X_n-\mu|\ge\varepsilon\bigr)\le\frac{\sigma^2}{n\varepsilon^2}.
\]
The right-hand side tends to zero for each fixed $\varepsilon>0$.
\end{proof}

Note what has just been proved and what has not. What is proved: for a tolerance
$\varepsilon$ that we choose ourselves, the probability that the average misses by
more than $\varepsilon$ tends to zero. What is not proved: how fast, and what the
shape of the remaining error is. Those are two different questions.

Figure~\ref{fig:tiga-lebar} draws the content of that statement for dice. Only two
things change from $n=4$ to $n=64$: the width shrinks and the peak rises, while the
center stays at $3.5$. The law of large numbers is not a promise that one run is
more accurate, but that the distribution of possibilities contracts to a single
point.

\begin{figure}[htbp]
\centering
\begin{tikzpicture}
\begin{axis}[
  width=0.95\linewidth, height=5.4cm,
  xmin=1, xmax=6, ymin=0, ymax=2.05,
  xlabel={value of the average $\bar x_n$}, ylabel={density},
  label style={font=\small}, tick label style={font=\footnotesize},
  axis lines=left,
  legend style={font=\footnotesize, draw=layergrey!35, fill=white,
    fill opacity=0.9, text opacity=1, at={(0.015,0.985)},
    anchor=north west, cells={anchor=west}},
  legend cell align=left,
]
\addplot[draw=layerteal!60!black, fill=layerteal!18, fill opacity=0.85,
  line width=0.4pt] coordinates{
  (1,0.0031) (1.25,0.0123) (1.5,0.0309) (1.75,0.0617) (2,0.108) (2.25,0.1728)
  (2.5,0.2469) (2.75,0.321) (3,0.3858) (3.25,0.4321) (3.5,0.4506) (3.75,0.4321)
  (4,0.3858) (4.25,0.321) (4.5,0.2469) (4.75,0.1728) (5,0.108) (5.25,0.0617)
  (5.5,0.0309) (5.75,0.0123) (6,0.0031)};
\addlegendentry{$n=4$}
\addplot[draw=layerindigo!70!black, fill=layerindigo!22, fill opacity=0.8,
  line width=0.4pt] coordinates{
  (1.6875,0) (1.75,0.0001) (1.8125,0.0002) (1.875,0.0004) (1.9375,0.0008) (2,0.0014)
  (2.0625,0.0025) (2.125,0.0043) (2.1875,0.0072) (2.25,0.0117) (2.3125,0.0184) (2.375,0.028)
  (2.4375,0.0416) (2.5,0.0601) (2.5625,0.0847) (2.625,0.1163) (2.6875,0.1558) (2.75,0.2039)
  (2.8125,0.2607) (2.875,0.3258) (2.9375,0.3981) (3,0.4759) (3.0625,0.5567) (3.125,0.6374)
  (3.1875,0.7145) (3.25,0.7843) (3.3125,0.8431) (3.375,0.8878) (3.4375,0.9156) (3.5,0.9251)
  (3.5625,0.9156) (3.625,0.8878) (3.6875,0.8431) (3.75,0.7843) (3.8125,0.7145) (3.875,0.6374)
  (3.9375,0.5567) (4,0.4759) (4.0625,0.3981) (4.125,0.3258) (4.1875,0.2607) (4.25,0.2039)
  (4.3125,0.1558) (4.375,0.1163) (4.4375,0.0847) (4.5,0.0601) (4.5625,0.0416) (4.625,0.028)
  (4.6875,0.0184) (4.75,0.0117) (4.8125,0.0072) (4.875,0.0043) (4.9375,0.0025) (5,0.0014)
  (5.0625,0.0008) (5.125,0.0004) (5.1875,0.0002) (5.25,0.0001) (5.3125,0)};
\addlegendentry{$n=16$}
\addplot[draw=accent!80!black, fill=accent!30, fill opacity=0.8,
  line width=0.5pt] coordinates{
  (2.5312,0) (2.5469,0.0001) (2.5625,0.0001) (2.5781,0.0001) (2.5938,0.0002) (2.6094,0.0003)
  (2.625,0.0004) (2.6406,0.0005) (2.6562,0.0007) (2.6719,0.0009) (2.6875,0.0012) (2.7031,0.0016)
  (2.7188,0.0021) (2.7344,0.0028) (2.75,0.0036) (2.7656,0.0047) (2.7812,0.0061) (2.7969,0.0079)
  (2.8125,0.01) (2.8281,0.0127) (2.8438,0.0161) (2.8594,0.0202) (2.875,0.0252) (2.8906,0.0312)
  (2.9062,0.0385) (2.9219,0.0472) (2.9375,0.0576) (2.9531,0.0698) (2.9688,0.0842) (2.9844,0.1009)
  (3,0.1203) (3.0156,0.1426) (3.0312,0.1681) (3.0469,0.1971) (3.0625,0.2297) (3.0781,0.2663)
  (3.0938,0.307) (3.1094,0.3521) (3.125,0.4015) (3.1406,0.4553) (3.1562,0.5136) (3.1719,0.5762)
  (3.1875,0.6429) (3.2031,0.7134) (3.2188,0.7875) (3.2344,0.8645) (3.25,0.944) (3.2656,1.0252)
  (3.2812,1.1075) (3.2969,1.19) (3.3125,1.2718) (3.3281,1.352) (3.3438,1.4295) (3.3594,1.5035)
  (3.375,1.573) (3.3906,1.6369) (3.4062,1.6943) (3.4219,1.7445) (3.4375,1.7867) (3.4531,1.8202)
  (3.4688,1.8445) (3.4844,1.8592) (3.5,1.8641) (3.5156,1.8592) (3.5312,1.8445) (3.5469,1.8202)
  (3.5625,1.7867) (3.5781,1.7445) (3.5938,1.6943) (3.6094,1.6369) (3.625,1.573) (3.6406,1.5035)
  (3.6562,1.4295) (3.6719,1.352) (3.6875,1.2718) (3.7031,1.19) (3.7188,1.1075) (3.7344,1.0252)
  (3.75,0.944) (3.7656,0.8645) (3.7812,0.7875) (3.7969,0.7134) (3.8125,0.6429) (3.8281,0.5762)
  (3.8438,0.5136) (3.8594,0.4553) (3.875,0.4015) (3.8906,0.3521) (3.9062,0.307) (3.9219,0.2663)
  (3.9375,0.2297) (3.9531,0.1971) (3.9688,0.1681) (3.9844,0.1426) (4,0.1203) (4.0156,0.1009)
  (4.0312,0.0842) (4.0469,0.0698) (4.0625,0.0576) (4.0781,0.0472) (4.0938,0.0385) (4.1094,0.0312)
  (4.125,0.0252) (4.1406,0.0202) (4.1562,0.0161) (4.1719,0.0127) (4.1875,0.01) (4.2031,0.0079)
  (4.2188,0.0061) (4.2344,0.0047) (4.25,0.0036) (4.2656,0.0028) (4.2812,0.0021) (4.2969,0.0016)
  (4.3125,0.0012) (4.3281,0.0009) (4.3438,0.0007) (4.3594,0.0005) (4.375,0.0004) (4.3906,0.0003)
  (4.4062,0.0002) (4.4219,0.0001) (4.4375,0.0001) (4.4531,0.0001) (4.4688,0)};
\addlegendentry{$n=64$}
\draw[layergrey, densely dotted, line width=0.6pt] (axis cs:3.5,0) -- (axis cs:3.5,1.95);
\node[font=\footnotesize, layergrey, anchor=south east] at (axis cs:3.45,1.95) {$3.5$};
\end{axis}
\end{tikzpicture}
\caption{The three distributions of the average of $n$ die rolls, drawn on the same
scale: $n=4$, $16$, and $64$. Their widths are $0.853913$, $0.426956$, and
$0.213478$, so each time $n$ is multiplied by four the width is halved. The peak
rises from $0.4506$ to $0.9251$ and $1.8641$, while the center stays at
$3.5$.}
\label{fig:tiga-lebar}
\end{figure}

The numbers make that difference tangible. For a coin toss with $p=1/2$ and
tolerance $\varepsilon=0.1$, Chebyshev gives
$\sigma^2/(n\varepsilon^2)=0.25/(0.01n)=25/n$. For the probability of missing to
be at most $0.3$, it suffices that $n\ge25/0.3=83.33$, that is $n=84$. At
$n=84$ that bound equals $0.297619$, so Chebyshev guarantees that the probability
of landing between $0.4$ and $0.6$ is at least $0.702381$. The true
probability, computed from the binomial distribution, is $0.937028$, and the
normal approximation gives $0.933197$. The guarantee is correct but too low by
about $23$ percentage points.

The same weakness becomes severe when $\sigma$ is large. For a fair die,
$\operatorname{Var}(X)=35/12=2.916667$ and $\sigma=1.707825$. At $n=100$ with
$\varepsilon=0.1$, Chebyshev gives the bound $2.916667$, a number larger than
one and therefore meaningless. To obtain an error of $0.01$ with $95$ percent
confidence, the normal formula demands at least $112043$ observations, whereas
Chebyshev demands far more still.

This is where the problem that sets up the next section lies. Chebyshev answers "may
the average be trusted" correctly but bluntly. It cannot answer "how close, with
what probability" in a usable way. A sharp answer demands knowledge of the shape of
the error distribution, not only of its size.

\begin{bacakeras}
Two voices, taking turns. The first names the number of tosses; the second names the
distance off that is still reasonable.

--- Ten tosses. The miss is still wide.
--- A hundred tosses. Already narrower.
--- A thousand tosses. Narrower still.

Now quadruple, do not add ten: ten becomes forty, a hundred becomes four hundred, a
thousand becomes four thousand. The miss is half. Not a quarter. Half. Every time
fourfold, exactly half.

Then reverse the direction, and read it more slowly. What is the cost of one more
digit after the decimal point? A hundred times the tosses. A hundred times the work,
one digit. The average is impatient: it demands four times, not two, and it offers
no discount.

The line repeated three times: four times the data, half the error.
\end{bacakeras}

\begin{untukanak}
Imagine a machine that fills bags of rice. Each bag misses the requested weight by a
little, sometimes over and sometimes under, and the size of the miss differs from
bag to bag; weighing one bag tells us nothing about the direction of its miss.

Now the machine fills a hundred bags, and we weigh all of them at once. The misses
upward and downward cancel each other, so the total weight is much closer to a
hundred times the requested weight. If there are ten thousand bags, the total is
almost exact.

The surprising part is that we do not need to know the shape of the distribution of
the miss. Whether the miss has a bell shape, takes only two values, or has any odd
shape at all, the distribution of the total still resembles the same bell, provided
there are enough bags and the miss does not turn into a giant too often. This is
what lets gamblers, bridge builders, and weather researchers use the same formula.

The question that makes us curious: why does the original shape of the miss no
longer matter once there are many bags? And why can one bag with a giant miss ruin
the whole conclusion? In this chapter we will find out why.
\end{untukanak}

\section{The Strong Law: What It Actually Adds}
\tingkat{advanced}

In everyday language: with confidence one, the average along one long run
eventually settles, rather than merely coming close often.

The weak law speaks about each $n$ separately: at a sufficiently large $n$, the
probability of missing is small. It does not forbid the average from missing badly
at $n=10^6$, then calming down, and then missing again at $n=10^9$, provided such
events are rare at each $n$. The strong law closes that gap.

\begin{teorema}[Strong law of large numbers]
Suppose $X_1,X_2,\dots$ are independent and identically distributed with
$\mathbb{E}|X_1|<\infty$ and $\mathbb{E}[X_1]=\mu$. Then
$\bar X_n\to\mu$ almost surely.\index{strong law of large numbers}
\end{teorema}

The difference is subtle but sharp. The sequence $X_n$ of the previous section
converges in probability to zero, yet at each $\omega$ its value is one infinitely
often. The weak law cannot reject it; the strong law does, because it demands
convergence at every $\omega$ outside a set of probability zero.

Note the hypothesis as well. The strong law demands only $\mathbb{E}|X_1|<\infty$,
with no variance. The weak law we proved above demands $\sigma^2<\infty$ only
because its proof uses Chebyshev; that hypothesis is not the natural hypothesis of
the theorem. For the weak law, a necessary and sufficient condition for an iid
sequence is $n\,P(|X_1|>n)\to0$, which is weaker than $\mathbb{E}|X_1|<\infty$. We
leave those details aside, and the reader can find them in Feller, \emph{An
Introduction to Probability Theory and Its Applications}, volume 2, in the chapter
on the law of large numbers.

We do not carry out the proof of the strong law here. Its classical form truncates
the variables at the level $n$, applies Borel--Cantelli to the tails of the
truncation, and then uses the Cauchy criterion for almost sure convergence; Etemadi
gives a proof that avoids the variance assumption and needs only pairwise
independence. Where the proofs live: Blitzstein and Hwang,
\emph{Introduction to Probability}, chapter 10, section 10.2; Wasserman, \emph{All
of Statistics}, appendix to chapter 5, section 5.7.1. What must be remembered from
that statement: the strong law gives no rate. It says that the error is eventually
small and stays small, not how fast.

\begin{tangan}
Take one coin and a sheet of lined paper. Heads counts one, tails counts zero.

Toss forty times, write 1 or 0 in a row. On the second line, write the running
average under each toss: the sum so far divided by the number of tosses so far.

What we should see: the second line swings most wildly over the first five tosses,
then grows calmer. At the fortieth toss the number need not be $0.5$; $0.475$ or
$0.525$ is still normal.

Now repeat that experiment four times and draw the four paths on one graph with
different pencils. All four paths narrow toward $0.5$, and the width of their
spread is halved each time the number of tosses is multiplied by four.

Tape the sheet to the wall. If someone says the results must balance, point to the
fortieth toss, not the tenth.
\end{tangan}

\section{The Central Limit Theorem}
\tingkat{advanced}

In everyday language: the error of the average, once magnified by a factor
$\sqrt n$, always has the same bell shape, whatever the original shape of the data.

The law of large numbers tells us that $\bar X_n$ sticks to $\mu$. It does not tell
us how wide that sticking is. The central limit theorem answers that, and the answer
is surprising because it does not depend on the original distribution.

\kamus{central limit theorem}{teorema limit pusat}{the bell shape that appears in
the error of the average, not in the data, and only after that error is magnified by
the square root of the number of data points}

\begin{teorema}[Central limit, Lindeberg--Levy form]\index{central limit theorem}
Suppose $X_1,X_2,\dots$ are independent and identically distributed with
$\mathbb{E}[X_i]=\mu$ and $0<\operatorname{Var}(X_i)=\sigma^2<\infty$. Then
\[
Z_n=\frac{\sqrt n\,(\bar X_n-\mu)}{\sigma}
\;\rightsquigarrow\; Z\sim N(0,1),
\]
meaning that $P(Z_n\le z)\to\Phi(z)$ for every $z\in\mathbb{R}$.
\end{teorema}

Note where $\sqrt n$ sits. The standard deviation of $\bar X_n$ is $\sigma/\sqrt n$,
not $\sigma/n$. That is why cutting the error in half requires four times the data,
and adding one digit of precision requires a hundred times the data. This is bad
news that cannot be dodged: precision comes at a quadratic price.

Figure~\ref{fig:lonceng-dari-kotak} shows where the bell shape comes from. The left
panel is the distribution of a single die: flat, with no bell at all. The right panel
is the distribution of the average of twelve dice after standardization. The raw
material is not a bell; the result is a bell.

\begin{figure}[htbp]
\centering
\begin{tikzpicture}
\begin{axis}[
  name=satu, width=4.3cm, height=4.5cm,
  ybar, bar width=11pt,
  xmin=0.5, xmax=6.5, ymin=0, ymax=0.21,
  xtick={1,2,3,4,5,6}, ytick={0,0.0833,0.1667},
  yticklabels={$0$,$\frac{1}{12}$,$\frac16$},
  xlabel={die face}, ylabel={probability},
  title={\footnotesize one die},
  label style={font=\footnotesize}, tick label style={font=\scriptsize},
  axis lines=left,
]
\addplot[fill=layerindigo!50, draw=layerindigo!80!black, line width=0.3pt]
  coordinates{
  (1,0.1667) (2,0.1667) (3,0.1667) (4,0.1667) (5,0.1667) (6,0.1667)};
\end{axis}
\begin{axis}[
  name=duabelas, at={(satu.east)}, anchor=west, xshift=7pt,
  width=9.3cm, height=4.5cm,
  ybar, bar width=3.4pt,
  xmin=-5.3, xmax=5.3, ymin=0, ymax=0.45,
  xtick={-4,-2,0,2,4}, ytick={0,0.2,0.4},
  xlabel={score $z$ from the mean}, ylabel={density},
  title={\footnotesize twelve dice, already standardized},
  label style={font=\footnotesize}, tick label style={font=\scriptsize},
  axis lines=left,
]
\addplot[fill=accent!45, draw=accent!70!black, line width=0.2pt]
  coordinates{
  (-4.3948,0) (-4.2258,0.00001) (-4.0567,0.00003) (-3.8877,0.00009) (-3.7187,0.0002) (-3.5496,0.00044)
  (-3.3806,0.00091) (-3.2116,0.00177) (-3.0426,0.00327) (-2.8735,0.00575) (-2.7045,0.00966) (-2.5355,0.01559)
  (-2.3664,0.02418) (-2.1974,0.03614) (-2.0284,0.05215) (-1.8593,0.07276) (-1.6903,0.09826) (-1.5213,0.1286)
  (-1.3522,0.16323) (-1.1832,0.20113) (-1.0142,0.24073) (-0.8452,0.28001) (-0.6761,0.31669) (-0.5071,0.34838)
  (-0.3381,0.37287) (-0.169,0.38835) (0,0.39365) (0.169,0.38835) (0.3381,0.37287) (0.5071,0.34838)
  (0.6761,0.31669) (0.8452,0.28001) (1.0142,0.24073) (1.1832,0.20113) (1.3522,0.16323) (1.5213,0.1286)
  (1.6903,0.09826) (1.8593,0.07276) (2.0284,0.05215) (2.1974,0.03614) (2.3664,0.02418) (2.5355,0.01559)
  (2.7045,0.00966) (2.8735,0.00575) (3.0426,0.00327) (3.2116,0.00177) (3.3806,0.00091) (3.5496,0.00044)
  (3.7187,0.0002) (3.8877,0.00009) (4.0567,0.00003) (4.2258,0.00001) (4.3948,0)};
\addplot[domain=-4.6:4.6, samples=120, layerindigo, line width=0.9pt]
  {0.39894228*exp(-x^2/2)};
\end{axis}
\end{tikzpicture}
\caption{From flat to bell. Left: the distribution of one die, flat over six values.
Right: the distribution of the average of twelve dice after standardization, stacked
with the standard normal curve. The supremum distance between the two is
$0.03384$, less than half the distance at $n=2$, which is $0.08333$. The raw
material is not a bell, the result is a bell.}
\label{fig:lonceng-dari-kotak}
\end{figure}

The following proof uses moment generating functions. It is complete for the case
in which the third moment is finite, and we state that limitation honestly at the
end.

\begin{lemma}[Convergence through moment generating functions]\index{moment generating function}
Suppose $\psi_n(t)=\mathbb{E}[e^{tZ_n}]$ and $\psi(t)=\mathbb{E}[e^{tZ}]$ are finite
on some open interval around $t=0$. If $\psi_n(t)\to\psi(t)$ for every $t$ in that
interval, then $Z_n\rightsquigarrow Z$.
\end{lemma}

That lemma is the continuity theorem for moment generating functions. Its proof uses
the fact that the distribution function is determined by its moment generating
function on an interval around zero, and we do not repeat it here; see Wasserman,
\emph{All of Statistics}, chapter 5, section 5.7.2, Lemma 5.20.

\begin{proof}[Proof of the central limit theorem]
Set $Y_i=(X_i-\mu)/\sigma$, so that $\mathbb{E}[Y_i]=0$,
$\operatorname{Var}(Y_i)=1$, and $Z_n=n^{-1/2}\sum_{i=1}^n Y_i$. Let
$\psi(t)=\mathbb{E}[e^{tY_1}]$ be the moment generating function of $Y_1$, which is
finite on the interval $|t|\le\delta$. Since the variables are independent, the
moment generating function of $Z_n$ is
\[
\xi_n(t)=\Bigl[\psi\bigl(t/\sqrt n\bigr)\Bigr]^n .
\]
Derivatives at zero give $\psi(0)=1$, $\psi'(0)=\mathbb{E}[Y_1]=0$, and
$\psi''(0)=\mathbb{E}[Y_1^2]=1$. The second-order Taylor expansion with third-order
remainder, valid when $\mathbb{E}|Y_1|^3<\infty$, reads
\[
\psi(s)=1+\frac{s^2}{2}+R(s),\qquad |R(s)|\le C|s|^3 \quad (|s|\le\delta).
\]
For $n$ large enough we have $|t|/\sqrt n\le\delta$, so
\[
\xi_n(t)=\Bigl(1+\frac{t^2}{2n}+R\bigl(t/\sqrt n\bigr)\Bigr)^n
=\Bigl(1+\frac{t^2}{2n}+O\bigl(n^{-3/2}\bigr)\Bigr)^n .
\]
Take logarithms. With $\log(1+a)=a-a^2/2+\dots$ and
$a_n=t^2/(2n)+O(n^{-3/2})$,
\[
\log\xi_n(t)=n\Bigl(a_n-\frac{a_n^2}{2}+\dots\Bigr)
=\frac{t^2}{2}+O\bigl(n^{-1/2}\bigr)-\frac{t^4}{8n}+\dots
\longrightarrow \frac{t^2}{2}.
\]
Hence $\xi_n(t)\to e^{t^2/2}$ for every fixed $t$. That last function is the moment
generating function of $N(0,1)$, so the lemma above gives
$Z_n\rightsquigarrow N(0,1)$.
\end{proof}

There is one further hypothesis that slips in: the third-order Taylor remainder
demands $\mathbb{E}|Y_1|^3<\infty$. For distributions with very heavy tails that
hypothesis can fail even though the variance is finite, and the proof above has to
be replaced. The usual route is to use the characteristic function
$\varphi(t)=\mathbb{E}[e^{itX}]$, which always exists because $|e^{itX}|=1$, and then
to apply Levy's continuity theorem. The skeleton of the proof is the same: the
logarithm of the characteristic function of $Z_n$ tends to $-t^2/2$. The version
without the third-moment assumption is in Wasserman, chapter 5, section 5.7.2, and
in Blitzstein and Hwang, chapter 10, section 10.3.

This theorem does not say that the error is small. It says that the error has a bell
shape once it is divided by $\sigma/\sqrt n$. Its size is still given by
$\sigma/\sqrt n$, and $\sigma$ is a quantity we usually do not know. In practice
$\sigma$ is replaced by the sample standard deviation $s$, and that replacement does
not destroy the limit; the statement becomes $\sqrt n(\bar X_n-\mu)/s_n
\rightsquigarrow N(0,1)$. The quantity $\sigma/\sqrt n$ is called the standard error
\index{standard error} (\emph{galat baku}), and it is the basis of the confidence
intervals in Chapter 20.

How fast that approximation improves can also be bounded.

\begin{teorema}[Berry--Esseen inequality]
If $\rho=\mathbb{E}|X_1-\mu|^3<\infty$, then
\[
\sup_z\bigl|P(Z_n\le z)-\Phi(z)\bigr|\le \frac{C\rho}{\sigma^3\sqrt n},
\]
where $C$ is a universal constant; the value $C=0.4748$ suffices for all
distributions.
\end{teorema}

The shape of this inequality matters: the distance shrinks like $1/\sqrt n$, with a
constant that depends on the skewness of the original distribution through
$\rho/\sigma^3$. A strongly skewed distribution will need a far larger $n$ than a
symmetric one. The numbers in the lab section show this quantitatively for the
exponential distribution, whose third moment is $\rho=2.414553$ and whose bound is
therefore $1.1464/\sqrt n$.

\begin{resep}
Three steps to standardize a sample average. The order must not be swapped.
\begin{enumerate}
\item \textbf{Subtract the center.} Compute $\bar x-\mu$. The bell shifts so that
      its peak falls at zero.
\item \textbf{Divide by the width.} Divide the result by the standard error
      $\sigma/\sqrt n$. The shape of the bell does not change; only its scale becomes
      one unit.
\item \textbf{Read the table.} The remaining number $z$ is read off $\Phi(z)$; for
      the right tail use $1-\Phi(z)$.
\end{enumerate}
A warning about order: dividing before subtracting the center shifts the result by
$\mu\sqrt n/\sigma$. For dice with $\mu=3.5$, $\sigma=1.707825$, and $n=12$,
that shift is $7.10$; the number is already outside any normal table.
\end{resep}

\section{The Normal Approximation and the Continuity Correction}
\tingkat{intermediate}

In everyday language: to compute the probability of a sum we may use the normal
table, but when the sum can only take integer values, the threshold is shifted by
one half.

The normal approximation for sums works most simply on discrete distributions, and
there a small correction appears whose effect is large. It is called the continuity
correction\index{continuity correction}. Take $S\sim
\text{Binomial}(100;\,0.5)$, so that $\mu=50$ and $\sigma=\sqrt{100\cdot0.5
\cdot0.5}=5$. The event $\{S\le40\}$ is the same event as $\{S<40.5\}$, because
$S$ only takes integer values. Hence the correct threshold to standardize is
$40.5$, not $40$.

\begin{center}
\begin{tabular}{lrrr}
\hline
Event & Exact & No correction & With correction \\
\hline
$P(S\le40)$ & $0.028444$ & $0.022750$ & $0.028717$ \\
$P(S\le55)$ & $0.864373$ & $0.841345$ & $0.864334$ \\
$P(S=60)$   & $0.010844$ & $0.010798$ & $0.010852$ \\
\hline
\end{tabular}
\end{center}

Read that table line by line. In the first row, the approximation without the
correction misses by $0.005694$, whereas with the correction it misses by
$0.000273$, almost twenty-one times closer. In the third row both are already
close, because near the center of the distribution a width of one unit is already
small compared with $\sigma=5$. The general lesson: the continuity correction helps
most in the tails and when $\sigma$ is small. The size of the shift in $z$ that it
produces is $1/(2\sigma)$, namely $0.1$ in this example.

\kertas{Recompute the first row of this table by hand. $z=(40-50)/5=-2$,
$\Phi(-2)=0.022750$; then $z=(40.5-50)/5=-1.9$, $\Phi(-1.9)=0.028717$.
The exact value is $\sum_{k=0}^{40}\binom{100}{k}2^{-100}=0.028444$. Check also
that the error without the correction $0.005694$ is about $21$ times the error
with the correction $0.000273$.}

This correction is not a new formula, only a rereading of the event. For a discrete
variable with spacing $h$ between values, the threshold is shifted by $h/2$ inward.
For a continuous variable nothing needs shifting, because there are no jumps.

One bad habit deserves mention here. Many books write that the normal approximation
"may be used when $n\ge30$". That is not a theorem but a habit. The numbers in the
lab section show why: for the exponential average, the supremum distance to the
normal at $n=30$ is $0.02429$, and the $95$ percent quantile of $Z_{30}$ is
$1.7419$ whereas the normal gives $1.6449$. That difference of $0.097$ means a
confidence interval that is too narrow by about six percent. What decides is not $n$
alone, but the skewness and kurtosis of the original distribution.

\section{When the Central Limit Theorem Misleads: Heavy Tails}
\tingkat{advanced}

In everyday language: when one observation can occasionally be far larger than all
the others, the average no longer settles, and the bell does not appear.

\kamus{heavy tail}{ekor berat}{a probability of huge values that does not shrink
fast enough, so that one observation can outweigh all the rest}

All the theorems above use one hypothesis that is easy to pass over: finite variance.
The Cauchy distribution\index{Cauchy distribution}, with density
\[
f(x)=\frac{1}{\pi(1+x^2)},\qquad x\in\mathbb{R},
\]
does not have it. Its tail shrinks like $1/x$, which is too slow. The probability of
exceeding ten is $1-(2/\pi)\arctan 10=0.063451$, and the probability of exceeding
one hundred is still $0.006366$. Compare the standard normal, whose probability of
exceeding ten is practically zero. Even $\mathbb{E}|X|$ is not finite: the
expectation truncated at level $M$ equals $(2/\pi)\ln(1+M^2)$, namely $2.9381$ for
$M=10$, $8.7952$ for $M=10^3$, and $17.5905$ for $M=10^6$. That number grows
like a logarithm, so there is no limit.

The consequence is blunt, and not merely "the CLT is less accurate".

\begin{proposisi}
If $X_1,\dots,X_n$ are iid Cauchy, then $\bar X_n$ has a Cauchy distribution for
every $n$.
\end{proposisi}

\begin{proof}[Proof]
The characteristic function of the Cauchy distribution is $\varphi(t)=e^{-|t|}$.
Since the variables are independent, the characteristic function of $\bar X_n$ is
$\bigl[e^{-|t|/n}\bigr]^n=e^{-|t|}$, which is the Cauchy characteristic function
again. Since the characteristic function determines the distribution, $\bar X_n$ has
a Cauchy distribution.
\end{proof}

So the average of $10^5$ Cauchy observations is no more informative than one
observation. The law of large numbers fails, the central limit theorem fails, and
"standard error" means nothing. The numbers in the lab section show all three
symptoms at once. The median of $|\bar X_n|$ stays around one for all $n$
($1.004$, $1.008$, $1.040$, $0.973$), whereas for a distribution with finite
variance it must shrink like $1/\sqrt n$. The probability $|\bar X_n|>10$ stays
around $0.0635$, its theoretical value, for all $n$. The empirical standard
deviation jumps about: $33.04$, $863.61$, $23.87$, $379.53$ for $n=10$ up to
$10^4$. Those jumps are not a bug in the program. The sample standard deviation of
Cauchy data is itself a random variable with no typical value.

Figure~\ref{fig:cauchy-berjalan} shows the difference on one run. The running
average of uniform data sticks to $0.5$ after a few dozen draws. The running
average of Cauchy data jumps from $-0.0866$ to $4.4572$ at the twentieth draw,
because one value $90.79$ happened to arrive, and then wanders again. One giant
value is enough to cancel the entire history before it.

\begin{figure}[htbp]
\centering
\begin{tikzpicture}
\begin{axis}[
  width=0.95\linewidth, height=6.0cm,
  xmode=log, xmin=1, xmax=4000, ymin=-1.1, ymax=4.9,
  xlabel={number of draws $n$ (logarithmic scale)}, ylabel={running average},
  label style={font=\small}, tick label style={font=\footnotesize},
  axis lines=left,
  legend style={font=\footnotesize, draw=layergrey!35, fill=white,
    fill opacity=0.9, text opacity=1, at={(0.015,0.985)},
    anchor=north west, cells={anchor=west}},
  legend cell align=left,
]
\addplot[accent, line width=0.7pt] coordinates{
  (1,-0.253) (2,1.045) (3,0.582) (4,0.564) (5,0.553) (6,1.064)
  (7,-0.579) (8,-0.517) (9,-0.409) (10,-0.325) (11,-0.409) (12,-0.361)
  (13,-0.344) (14,-0.262) (15,0.008) (16,-0.038) (17,0.017) (18,-0.08)
  (19,-0.087) (20,4.457) (21,4.29) (22,4.153) (23,4.046) (24,3.762)
  (25,3.656) (26,3.551) (27,3.425) (28,3.257) (29,3.186) (30,3.083)
  (31,2.934) (32,2.786) (33,2.837) (34,2.743) (35,2.633) (36,2.528)
  (37,2.464) (38,2.381) (39,2.279) (40,2.218) (41,2.157) (44,1.732)
  (47,1.387) (50,1.367) (53,1.191) (57,1.113) (61,1.077) (65,0.942)
  (70,0.796) (75,0.773) (80,0.725) (85,0.492) (91,0.79) (97,0.655)
  (104,0.54) (111,0.478) (119,0.494) (127,0.401) (135,0.981) (145,0.868)
  (155,2.5) (165,2.351) (177,2.207) (189,2.168) (202,1.935) (216,1.854)
  (230,1.776) (246,1.665) (263,1.573) (281,0.848) (300,0.794) (321,0.791)
  (343,0.511) (367,0.922) (392,0.986) (419,0.984) (447,1.03) (478,0.951)
  (511,0.827) (546,0.685) (583,0.37) (623,0.548) (666,0.953) (712,1.021)
  (761,0.945) (813,2.066) (869,2.135) (929,1.855) (992,1.814) (1060,1.649)
  (1133,1.31) (1211,0.557) (1294,0.422) (1383,0.434) (1478,0.466) (1579,0.324)
  (1688,0.406) (1803,0.386) (1927,0.444) (2059,0.403) (2201,0.585) (2352,0.419)
  (2513,0.528) (2686,0.397) (2870,0.356) (3067,0.326) (3278,0.276) (3503,-0.096)
  (3743,-0.159) (4000,-0.203)};
\addlegendentry{Cauchy}
\addplot[layerteal, dashed, line width=1.0pt] coordinates{
  (1,0.221) (2,0.272) (3,0.323) (4,0.472) (5,0.462) (6,0.519)
  (7,0.524) (8,0.573) (9,0.55) (10,0.574) (11,0.605) (12,0.601)
  (13,0.57) (14,0.581) (15,0.543) (16,0.532) (17,0.529) (18,0.528)
  (19,0.532) (20,0.544) (21,0.545) (22,0.547) (23,0.539) (24,0.531)
  (25,0.54) (26,0.555) (27,0.561) (28,0.541) (29,0.545) (30,0.547)
  (31,0.562) (32,0.549) (33,0.533) (34,0.542) (35,0.531) (36,0.535)
  (37,0.537) (38,0.526) (39,0.519) (40,0.516) (41,0.517) (44,0.518)
  (47,0.519) (50,0.503) (53,0.518) (57,0.515) (61,0.525) (65,0.53)
  (70,0.529) (75,0.53) (80,0.535) (85,0.538) (91,0.539) (97,0.538)
  (104,0.541) (111,0.535) (119,0.532) (127,0.533) (135,0.528) (145,0.526)
  (155,0.518) (165,0.523) (177,0.513) (189,0.516) (202,0.513) (216,0.507)
  (230,0.5) (246,0.498) (263,0.495) (281,0.496) (300,0.496) (321,0.495)
  (343,0.499) (367,0.499) (392,0.501) (419,0.503) (447,0.497) (478,0.505)
  (511,0.508) (546,0.509) (583,0.504) (623,0.502) (666,0.497) (712,0.494)
  (761,0.488) (813,0.484) (869,0.485) (929,0.484) (992,0.485) (1060,0.484)
  (1133,0.483) (1211,0.485) (1294,0.487) (1383,0.489) (1478,0.488) (1579,0.491)
  (1688,0.489) (1803,0.489) (1927,0.49) (2059,0.491) (2201,0.49) (2352,0.492)
  (2513,0.491) (2686,0.493) (2870,0.492) (3067,0.49) (3278,0.492) (3503,0.491)
  (3743,0.493) (4000,0.493)};
\addlegendentry{uniform $(0,1)$}
\draw[layergrey, densely dotted, line width=0.6pt]
  (axis cs:1,0.5) -- (axis cs:4000,0.5);
\node[font=\scriptsize, layergrey, anchor=south west] at (axis cs:1.3,0.53) {$0.5$};
\end{axis}
\end{tikzpicture}
\caption{Running averages over 4000 draws. Dashed line: uniform on $(0,1)$, settling
near $0.5$ and ending at $0.4930$. Solid line: Cauchy, jumping from $-0.0866$
to $4.4572$ at the twentieth draw when one value $90.79$ arrived, then wandering
again and ending at $-0.203$. The horizontal axis is logarithmic in width, so that
the early jump is not hidden.}
\label{fig:cauchy-berjalan}
\end{figure}

What replaces the central limit theorem is the family of stable
distributions\index{stable distribution}. For iid variables with tails that shrink
like $x^{-\alpha}$ with $0<\alpha<2$, there are constants $a_n$ such that
\[
\frac{S_n-a_n}{n^{1/\alpha}}\;\rightsquigarrow\;Y_\alpha,
\]
with $Y_\alpha$ the stable distribution of index $\alpha$. For $\alpha=1$ we get
Cauchy and $n^{1/1}=n$, exactly as in the proposition above. For
$\alpha\in(1,2)$ the average still converges to $\mu$, but the error shrinks like
$n^{1-1/\alpha}$, slower than $1/\sqrt n$, and its shape is not a bell. This general
statement is the result of Gnedenko and Kolmogorov; the details are in Feller volume
2, in the chapter on stable distributions.

If $\alpha>2$ but the third moment does not exist, the central limit theorem still
holds because the variance is finite. What is lost is the guarantee of a rate: the
Berry--Esseen bound becomes infinite, so there is no longer any handle on the claim
that the error shrinks like $1/\sqrt n$ with a controlled constant. In practice this
is what bites most often: the distribution is not extreme like Cauchy, but it has
enough tail to make confidence intervals too narrow at realistic sample sizes.

\begin{tantangan}
A distribution with no mean. Take $X$ with
$P(X=2^k)=P(X=-2^k)=2^{-k-1}$ for $k=1,2,3,\dots$

Show two things. First, $\mathbb{E}|X|=\infty$, so there is no $\mu$ that may be
approached. Second, the sample average does not converge in probability to any
constant: the law of large numbers fails, it does not merely slow down.

Hint 1: what decides is not the average of all the terms, but the largest term. Ask
how often one term exceeds the threshold $n$.

Hint 2: compute $n\,P(|X|>n)$ for $n$ between two powers of two, then use the
necessary condition $n\,P(|X_1|>n)\to0$ stated in the section on the strong law.

If you get stuck, go back to the section on heavy tails and to the paragraph on the
necessary and sufficient condition in the section on the strong law.
\end{tantangan}

\begin{untukanak}
In a village there are a hundred people with ordinary incomes, and one person who is
extremely rich. If we compute the average income of that village, the number is
decided almost entirely by that rich person. If the rich person happens to be away,
the average drops drastically. If that person comes back, the average jumps up.

Now add inhabitants. Another hundred people, another thousand; as long as the
newcomers are ordinary people, the average changes little. But if among the newcomers
there is one more giant, the average jumps again. Adding many people does not help,
because what decides is not the number of people but who is the largest.

This is what happens with data whose tails are heavy: one observation can outweigh
all the others. The average becomes unstable, and the bell that is usually used to
compute "how confident" does not appear at all.

The question that makes us curious: how can we know, before computing, whether our
data are of the calm kind or of the kind ruled by giants? In this chapter we will
find out why that question cannot be answered just by looking at the average.
\end{untukanak}

\section{Computational Lab: Five Checks}
\tingkat{advanced}

The five claims of this chapter can be checked with one script: the $1/\sqrt n$
rate in the law of large numbers, the skewness and kurtosis of the exponential
average, the supremum distance to the normal, the Monte Carlo standard error, and
total failure for Cauchy. All the numbers below are real output from an execution
with NumPy 2.5.2 and SciPy 1.18.0, generator \texttt{default\_rng(20261024)}.

\begin{lab}
\begin{verbatim}
import numpy as np
from scipy import stats

rng = np.random.default_rng(20261024)

print("== 1. law of large numbers: dice average ==")
for k in range(1, 6):
    n = 10**k
    m = rng.integers(1, 7, size=(2000, n)).mean(axis=1)
    print("n=%6d  mean=%.4f  rms=%.4f  sd/sqrt n=%.4f"
          % (n, m.mean(), np.sqrt(np.mean((m-3.5)**2)),
             np.sqrt(35/12/n)))

print("== 2. CLT: exponential average ==")
for n in (2, 5, 10, 30, 100, 500):
    z = (rng.exponential(size=(50000, n)).mean(axis=1)
         - 1)*np.sqrt(n)
    ks = stats.kstest(z, "norm")
    print("n=%4d skew=%+.3f (%+.3f) kurt=%+.3f (%+.3f)"
          % (n, stats.skew(z), 2/np.sqrt(n),
             stats.kurtosis(z), 6.0/n))
    print("      KS D=%.4f p=%.3f" % (ks.statistic, ks.pvalue))

print("== 3. sup distance to normal, exact via gamma ==")
rho = 2 + 2*(6/np.e - 2)          # E|X-1|^3 for X ~ Exp(1)
for n in (2, 5, 10, 30, 100, 500):
    zg = np.linspace(-4, 4, 40001)
    Fn = stats.gamma.cdf(n*(1+zg/np.sqrt(n)), a=n)
    D = np.max(np.abs(Fn - stats.norm.cdf(zg)))
    print("n=%4d sup|Fn-Phi|=%.5f  D*sqrt n=%.3f  BE=%.3f"
          % (n, D, D*np.sqrt(n), 0.4748*rho/np.sqrt(n)))

print("== 4. Monte Carlo standard error for E[exp(U)] ==")
for k in (2, 3, 4, 5):
    n = 10**k
    est = np.exp(rng.random(size=(2000, n))).mean(axis=1)
    print("n=%6d  sd=%.5f  sd*sqrt n=%.5f"
          % (n, est.std(), est.std()*np.sqrt(n)))

print("== 5. Cauchy: the average does not concentrate ==")
for k in (1, 2, 3, 4):
    n = 10**k
    m = rng.standard_cauchy(size=(4000, n)).mean(axis=1)
    print("n=%5d  median|m|=%.3f P(|m|>10)=%.4f sd=%.2f"
          % (n, np.median(np.abs(m)), np.mean(np.abs(m) > 10),
             m.std()))
print("theory P(|mean|>10) = %.4f"
      % (1 - 2/np.pi*np.arctan(10)))
\end{verbatim}
\end{lab}

Its output, verbatim:

\begin{verbatim}
== 1. law of large numbers: dice average ==
n=    10  mean=3.5014  rms=0.5440  sd/sqrt n=0.5401
n=   100  mean=3.4992  rms=0.1698  sd/sqrt n=0.1708
n=  1000  mean=3.4986  rms=0.0527  sd/sqrt n=0.0540
n= 10000  mean=3.5001  rms=0.0171  sd/sqrt n=0.0171
n=100000  mean=3.4997  rms=0.0053  sd/sqrt n=0.0054
== 2. CLT: exponential average ==
n=   2 skew=+1.440 (+1.414) kurt=+3.130 (+3.000)
      KS D=0.0929 p=0.000
n=   5 skew=+0.919 (+0.894) kurt=+1.393 (+1.200)
      KS D=0.0602 p=0.000
n=  10 skew=+0.631 (+0.632) kurt=+0.589 (+0.600)
      KS D=0.0480 p=0.000
n=  30 skew=+0.355 (+0.365) kurt=+0.181 (+0.200)
      KS D=0.0219 p=0.000
n= 100 skew=+0.196 (+0.200) kurt=+0.038 (+0.060)
      KS D=0.0167 p=0.000
n= 500 skew=+0.112 (+0.089) kurt=-0.003 (+0.012)
      KS D=0.0062 p=0.043
== 3. sup distance to normal, exact via gamma ==
n=   2 sup|Fn-Phi|=0.09446  D*sqrt n=0.134  BE=0.811
n=   5 sup|Fn-Phi|=0.05963  D*sqrt n=0.133  BE=0.513
n=  10 sup|Fn-Phi|=0.04211  D*sqrt n=0.133  BE=0.363
n=  30 sup|Fn-Phi|=0.02429  D*sqrt n=0.133  BE=0.209
n= 100 sup|Fn-Phi|=0.01330  D*sqrt n=0.133  BE=0.115
n= 500 sup|Fn-Phi|=0.00595  D*sqrt n=0.133  BE=0.051
== 4. Monte Carlo standard error for E[exp(U)] ==
n=   100  sd=0.04928  sd*sqrt n=0.49281
n=  1000  sd=0.01586  sd*sqrt n=0.50156
n= 10000  sd=0.00495  sd*sqrt n=0.49538
n=100000  sd=0.00155  sd*sqrt n=0.49056
== 5. Cauchy: the average does not concentrate ==
n=   10  median|m|=1.004 P(|m|>10)=0.0630 sd=33.04
n=  100  median|m|=1.008 P(|m|>10)=0.0680 sd=863.61
n= 1000  median|m|=1.040 P(|m|>10)=0.0653 sd=23.87
n=10000  median|m|=0.973 P(|m|>10)=0.0580 sd=379.53
theory P(|mean|>10) = 0.0635
\end{verbatim}

Part 1 shows the law of large numbers at work, and shows its price. The rms column
matches $\sigma/\sqrt n$ to two decimal places for $n\ge100$: $0.1698$ against
$0.1708$, $0.0171$ against $0.0171$, $0.0053$ against $0.0054$. But look
at the magnitudes: at $n=10^5$ the average still misses by about $0.005$, which is
$0.15$ percent of $3.5$. The average of $10^5$ tosses is not "practically
exact"; it is exact to two decimal places.

Parts 2 and 3 show two ways of measuring closeness to the normal. Skewness and
kurtosis follow the theoretical prediction neatly from $n=10$ on: $0.631$ against
$2/\sqrt{10}=0.632$ and $0.589$ against $6/10=0.600$. At $n=2$ and $n=5$ the
kurtosis estimator is still noisy, because its deviation depends on the eighth moment
of that distribution, not only on the number of repetitions. Part 3 computes the
supremum distance exactly, without simulation, by using the Gamma distribution for
the exponential average: $0.04211$ at $n=10$ and $0.00595$ at $n=500$. The column
$D\sqrt n$ is nearly constant at $0.133$, direct evidence that the rate is
$1/\sqrt n$. The Berry--Esseen bound $1.1464/\sqrt n$ holds, with the numbers
$0.209$ against $0.02429$ at $n=30$; that bound is about $8.6$ times looser
than reality, which is normal for a bound that is uniform over all distributions.

The Kolmogorov--Smirnov test on fifty thousand draws rejects normality with
$p=0.000$ for all $n$ up to $100$, then gives $p=0.043$ at $n=500$. That is not a
sign that the program is wrong: the normal approximation is never exact for finite
$n$, and the distance merely shrinks until a test that large can no longer see it.
The right question is not "is it normal" but "how large an error can we still
accept", and part 3 answers it with numbers.

Part 4 separates two quantities that are often confused. The standard deviation of
the estimator falls from $0.04928$ to $0.00155$ when $n$ grows a thousandfold,
whereas $\text{sd}\cdot\sqrt n$ is nearly constant at $0.49$. That constant value
is the $\sigma$ of $\exp(U)$, namely $\sqrt{(e^2-1)/2-(e-1)^2}=0.491971$. This is
the $1/\sqrt n$ law in its barest form: for a precision of $0.001$ at $95$ percent
confidence we need $(1.96\cdot0.491971/0.001)^2=929803$ draws.

Part 5 is the counterweight. The median of $|\bar X_n|$ does not shrink; the
probability of exceeding ten does not shrink; and the empirical standard deviation
jumps from $33.04$ to $863.61$ and then down to $23.87$. Theory says
$0.0635$, and all five simulation numbers hover around it. There is no law of large
numbers here, and no central limit theorem. The only way to rescue the analysis is to
change the question: estimate the tail index, not the mean.

\begin{untukdosen}
The teaching order I suggest: start from the simulation table, not from the
definitions. Show the class that the error of one run follows no pattern at all,
whereas the distribution of the error does. After that, the three modes of
convergence, with a counterexample for each direction that fails; do not postpone
those examples, because that is where students understand that "approaching" needs to
be made precise. Only then Markov, Chebyshev, the weak law, the strong law as a
statement, and then the central limit theorem with the moment generating function
proof.

Two misconceptions appear most often. First, students read "converges in
probability" as "the value becomes fixed", and are then surprised when a simulation
shows big jumps at large $n$; the split $(0,1)$ example in this chapter is the
antidote. Second, they read the central limit theorem as "the average becomes
normal", whereas what becomes normal is the error divided by $\sigma/\sqrt n$; as a
result they write $\bar X_n\approx N(\mu,\sigma^2)$ and forget $\sqrt n$.

A blackboard activity that works: ask two students to play the gambler and the
bookmaker. The gambler names a tolerance, the bookmaker names the number of tosses
that guarantees it, and the class computes with Chebyshev and then compares against
the exact numbers from SciPy. The gap is always large, and that is the doorway into
the normal approximation. Sign that the class is ready to move on: they can explain
why Chebyshev is still useful although it is loose, and they can name one
distribution that makes the central limit theorem fail.
\end{untukdosen}

\section{Summary}

\begin{ringkasan}
\begin{itemize}
  \item There are three different modes of convergence: almost sure, in probability,
        and in distribution. Almost sure implies in probability, in probability
        implies in distribution, and the reverse directions fail unless the limit is
        a constant.
  \item Markov's and Chebyshev's inequalities use only moments, which is why they
        hold for every distribution with finite variance, and why they are loose.
  \item Weak law of large numbers: $P(|\bar X_n-\mu|\ge\varepsilon)\le
        \sigma^2/(n\varepsilon^2)\to0$. The strong law adds that
        $\bar X_n\to\mu$ almost surely when $\mathbb{E}|X_1|<\infty$, without
        demanding a variance.
  \item Central limit theorem: $\sqrt n(\bar X_n-\mu)/\sigma\rightsquigarrow N(0,1)$
        when the variance is finite. The rate is $1/\sqrt n$, hence four times the
        data for half the error, and a hundred times the data for one extra digit.
  \item The normal approximation for discrete sums requires a continuity correction
        of half a unit; in the binomial example of this chapter that correction cuts
        the error in the tail by about a factor of twenty.
  \item If the variance is infinite, the central limit theorem fails, it does not
        merely slow down. The average of $n$ Cauchy observations is still Cauchy, and
        what replaces the bell is a stable distribution of index $\alpha<2$.
\end{itemize}
\end{ringkasan}

\section{Exercises}

\begin{latihan}
\begin{enumerate}
  \item Prove: if $X_n\to X$ and $Y_n\to Y$ in probability, then
        $X_n+Y_n\to X+Y$ in probability. Write down the set inequality you use, then
        explain why $\varepsilon/2$ appears.
  \item Prove that Chebyshev's inequality is sharp: for $0<\sigma^2\le a^2$, take $X$
        with $P(X=a)=P(X=-a)=\sigma^2/(2a^2)$ and
        $P(X=0)=1-\sigma^2/a^2$, then compute $\operatorname{Var}(X)$ and
        $P(|X|\ge a)$.
  \item Evaluate this claim: "Since the central limit theorem holds for all
        distributions, the average of thirty observations may always be treated as
        normal." True, false, or not enough evidence? Name the hypothesis that is
        violated or that has not been checked, and use numbers from the lab section.
  \item For a coin with $p=1/2$ and tolerance $\varepsilon=0.05$, find the smallest
        $n$ that Chebyshev guarantees for a probability of missing of at most
        $0.05$. Then find the smallest $n$ that really satisfies that condition, by
        computing the binomial distribution in Python. Compare the two numbers.
  \item Prove that convergence in probability does not guarantee convergence of
        expectations, using $P(X_n=n^2)=1/n$ and $P(X_n=0)=1-1/n$. Compute
        $\mathbb{E}[X_n]$ and $P(|X_n|>\varepsilon)$ for fixed $\varepsilon$.
  \item For $S\sim\text{Binomial}(200;\,0.3)$, approximate $P(S\le50)$ with the
        normal without and with the continuity correction, then compare both with the
        exact value from SciPy. Report both errors.
  \item Generate $4000$ averages of $n$ Cauchy numbers for
        $n=10,10^2,10^3,10^4$, then report the median of $|\bar X_n|$, the
        probability $|\bar X_n|>10$, and the standard deviation. Compare with
        $1-(2/\pi)\arctan10$ and explain why the standard deviation column does not
        shrink.
  \item Prove with characteristic functions that the average of $n$ iid Cauchy
        variables has a Cauchy distribution. Then explain why a report of the form
        "average $\pm$ two standard errors" on such data cannot be interpreted.
  \item How many Monte Carlo draws are needed for the estimator
        $\mathbb{E}[\exp(U)]$ to have a $95$ percent interval half-width of
        $0.0005$? Use $\sigma=0.491971$ from the lab section, then compare with
        the $n$ needed for a half-width of $0.001$. State the ratio and explain
        where it comes from.
  \item Prove that if $X_n\rightsquigarrow c$ with $c$ a constant, then
        $X_n\to c$ in probability, without using part (c) of the theorem on the
        relations among the three modes at the start of the chapter, then show by
        example that the hypothesis "constant" cannot be dropped.
\end{enumerate}

Hints for the three hardest items. Item 1: write
$\{|(X_n+Y_n)-(X+Y)|>\varepsilon\}$ as contained in the union of two events, then use
the union bound. Item 8: the characteristic function of a sum is the product of the
characteristic functions, and the Cauchy characteristic function $e^{-|t|}$ is
special because its form does not change when it is raised to a power. Item 10: write
$P(X_n\le t)\le P(X_n\le c-\varepsilon)+P(|X_n-c|>\varepsilon)$ for $t$ to the left
of $c$, and remember that a point-mass distribution jumps exactly at $c$.
\end{latihan}

\begin{untukpeneliti}
What is really proved in this chapter: the implications almost sure $\Rightarrow$ in
probability $\Rightarrow$ in distribution, the weak law via Chebyshev, and the
central limit theorem via moment generating functions under the assumption
$\mathbb{E}|Y_1|^3<\infty$. What is popular but not proved here: the strong law (the
statement only, Etemadi's proof is skipped), the version of the central limit theorem
without a third moment (it needs characteristic functions and Levy's continuity
theorem), and the whole theory of stable distributions (the Gnedenko--Kolmogorov
statement, without proof).

Primary sources I used: Wasserman, \emph{All of Statistics}, chapter 5, sections 5.2
(types of convergence and the implication diagram), 5.3 (the weak law), 5.4 (the
central limit theorem and Berry--Esseen), plus appendix 5.7.1 and 5.7.2 for the strong
law and the proof of the central limit theorem. Blitzstein and Hwang,
\emph{Introduction to Probability}, chapter 10, sections 10.1 (inequalities), 10.2
(the law of large numbers), 10.3 (the central limit theorem). For the heavy-tail
side: Feller volume 2.

Limits of validity that must be mentioned. First, all the theorems here assume
independence and identical distribution; for dependent data, for instance time series
or iterates of a Markov chain, what holds are martingale or mixing versions, and the
rate can be far slower. Second, the optimal constant in Berry--Esseen is not known;
$0.4748$ suffices, and the best lower bound for that constant is still below
$0.5$. Third, the rate $1/\sqrt n$ is the worst rate that can be guaranteed
uniformly; for particular distributions the true rate is better, and for heavy tails
with $\alpha\in(1,2)$ the true rate is worse. Open questions that are still alive:
how sharp the Berry--Esseen bound is for classes of distributions with bounded
moments, and how to estimate the tail index with valid confidence intervals at small
$n$.
\end{untukpeneliti}

\peta{Chapter 19 in Wasserman, \emph{All of Statistics}: chapter 5 (Convergence of
Random Variables), sections 5.2 Types of Convergence, 5.3 The Law of Large Numbers,
5.4 The Central Limit Theorem, and appendix 5.7.1 Almost Sure and L1 Convergence as
well as 5.7.2 Proof of the Central Limit Theorem. In Blitzstein and Hwang,
\emph{Introduction to Probability}: chapter 10 (Inequalities and limit theorems),
sections 10.1 Inequalities, 10.2 Law of large numbers, 10.3 Central limit theorem.
For stable distributions: Feller, \emph{An Introduction to Probability Theory and Its
Applications}, volume 2. The basic distributions used here are built in Chapter 17,
and the confidence intervals that rest on this chapter are discussed in Chapter 19.}
